\documentclass[12pt]{report}

\usepackage[utf8]{inputenc}   
\usepackage[top=2cm, bottom=2cm, left=2cm, right=2cm]{geometry}
\usepackage{color}
\usepackage{array}
\usepackage[T1]{fontenc}
\usepackage{hyperref}
\usepackage{graphicx}
\usepackage { boxedminipage }


\begin{document}
\begin{minipage}{0.55\linewidth}
\noindent Assocation : DaVinciBot\\
Ecole : ESILV - Léonard de vinci\\
Siège :  12 Avenue Léonard de Vinci, 92400 Courbevoie\\
Equipe Projet ESILV : CFR 2017\\
\end{minipage}\hfill
\begin{minipage}{0.3\linewidth}
\includegraphics[scale = 0.2]{img/logo_assos.png}
\end{minipage}
\vspace{1cm}
% PAS TOUCHER AU DESSUS !!! 


% modifier numéro compte rendu et numéro de semaine

\begin{center}
\begin {boxedminipage} [ poslb ] { 8cm}
\begin{center}
\textbf{Compte Rendu n01}\\
\textbf{Date 10 septemebre 2016}
\end{center}
\end {boxedminipage}
\end{center}
\textbf{Chargé du CR : Hugo Pousseur}

\vspace{1cm}

\paragraph{Etaient présent à la réunion :}
\begin{itemize}
\item Tout le monde
\end{itemize}


\paragraph{Ordre du jour :}
\begin{enumerate}
 \item Explication de la CFR 
 \item Choix de la méthode, SCRUMP Implique une dévision du projet en petit projet, en sprint et *zone de crise*.
 \item Explication du système de repérage.(But principale de cette année)
 \item Discussions sur des balises (Avec M Peretti)
 \item Discussions sur la différene entre signal sinusoidale et carré (Avec M Peretti)
\end{enumerate}


\paragraph{Décisions prises :}
\begin{enumerate}
 \item Taches définies
 \item Insallation de Kinetic (ROS)
 \item Etablir le planning le diagramme Pert
 \item Création d'un IRC (freenode) CFR\_ESILV
 \item Tous nos documents en LATEX
 \item Documentation à la fin d'une tâche en LATEX
 \item Sponsor
 \item Linux (Ubuntu)
 \item ROS (Kinetic)
\end{enumerate}


\paragraph{Ordre du jour de la prochaine réunion :}
\begin{enumerate}
 \item Réviser Python/Ros 
\end{enumerate}




\end{document}